%% ma: L12-B04-0016
%% lop: 12
%% chuong: 1
%% bai: 4
%% dang: 12.4.9
%% loai: TN4
%% muc: TH
%% dap_an: "A"
%% nguon: "HK1-DeOn_1_1.pdf (M10, Đề ôn tập học kì 1 Toán 12, Đề số 2), Phần I Câu 10"
%% kien_thuc: "Bài 3 (đường tiệm cận), Bài 4 (khảo sát và vẽ đồ thị hàm số)"
%% trang_thai: chinh-thuc
\begin{ex}
Tâm đối xứng của đồ thị hàm số $y=x-1+\dfrac{4}{x+2}$ có toạ độ là
\choice{\True $(-2;-3)$}{$(2;-3)$}{$(-2;3)$}{$(2;3)$}
\loigiai{Đồ thị có tiệm cận đứng $x=-2$ và tiệm cận xiên $y=x-1$; tâm đối xứng là giao điểm hai tiệm cận: $(-2;-3)$.}
\end{ex}
